\documentclass[a4paper,11pt]{article}

\usepackage[T1]{fontenc}
\usepackage[utf8]{inputenc}
\usepackage{lmodern}
\usepackage{enumitem}
\usepackage[hidelinks]{hyperref}
\usepackage{geometry}
\usepackage{amsmath}

\geometry{margin=2.5cm}
\setlist[itemize]{leftmargin=1.2cm}

\begin{document}

\begin{center}
    {\Large\bfseries Continuous interpretation of ARD}
\end{center}

\vspace{0.8cm}

\section{A continuous interpretation of ARD}

The purpose of this section is to interpret the ARD correction as the
discretization of an interface-supported distribution generated by the
artificial decomposition. To this end, we adopt the immersed-layer framework
introduced by Eldredge \cite{Eldredge2021}.

\subsection{Distributional formulation}

Let
\[
\Omega=\Omega^-\cup\Gamma\cup\Omega^+
\]
be the domain, where $\Gamma$ denotes the interface separating two neighboring
subdomains. We define a signed-distance function $\chi(x)$ such that
\[
\Gamma=\{x:\chi(x)=0\},
\]
with $\chi(x)<0$ in $\Omega^-$ and $\chi(x)>0$ in $\Omega^+$. The unit normal
vector
\[
n=\nabla\chi
\]
therefore points from $\Omega^-$ to $\Omega^+$.

Denote by $p^+$ and $p^-$ the pressure fields restricted to $\Omega^+$ and
$\Omega^-$, respectively. We introduce the piecewise field
\begin{equation}
\bar p
=
H(\chi)p^+
+
H(-\chi)p^-,
\label{eq:pbar}
\end{equation}
where $H$ denotes the Heaviside function. Similarly, we define
\begin{equation}
\overline{\Delta p}
=
H(\chi)\Delta p^+
+
H(-\chi)\Delta p^-.
\label{eq:piecewise-laplacian}
\end{equation}

Following Eldredge \cite[Eq.~(4)]{Eldredge2021}, differentiation of
\eqref{eq:pbar} yields
\begin{equation}
\Delta\bar p
=
\overline{\Delta p}
+
[\partial_n p]\delta_\Gamma
+
\nabla\cdot
\big(
[p]\delta_\Gamma n
\big),
\label{eq:eldredge}
\end{equation}
where
\[
[p]=p^+-p^-,
\qquad
[\partial_n p]
=
\partial_n p^+
-
\partial_n p^-,
\]
and $\delta_\Gamma=\delta(\chi)$ denotes the surface delta distribution
concentrated on $\Gamma$.

The two singular terms correspond respectively to a single-layer
(monopole) source and a double-layer (dipole) source. For a locally flat
interface, or equivalently in one dimension,
\[
\nabla\cdot([p]\delta_\Gamma n)
=
[p]\partial_n\delta_\Gamma,
\]
so that
\begin{equation}
\Delta\bar p
=
\overline{\Delta p}
+
[\partial_n p]\delta_\Gamma
+
[p]\partial_n\delta_\Gamma.
\label{eq:eldredge-flat}
\end{equation}

\subsection{Continuous interpretation of the ARD correction}

In the ARD method considered here, each subdomain is evolved independently
using homogeneous Neumann conditions on the artificial interfaces,
\[
\partial_n p^+
=
\partial_n p^-
=
0.
\]
Hence
\[
[\partial_n p]=0,
\]
while a jump in the pressure itself may develop. Equation
\eqref{eq:eldredge} therefore reduces to
\begin{equation}
\Delta\bar p
=
\overline{\Delta p}
+
\nabla\cdot([p]\delta_\Gamma n).
\label{eq:continuous-neumann}
\end{equation}

For a locally flat interface,
\[
\nabla\cdot([p]\delta_\Gamma n)
=
[p]\partial_n\delta_\Gamma,
\]
and thus
\begin{equation}
\Delta\bar p
=
\overline{\Delta p}
+
[p]\partial_n\delta_\Gamma.
\label{eq:continuous-neumann-flat}
\end{equation}

Consequently, the wave equation becomes
\begin{equation}
\bar p_{tt}
=
c^2\overline{\Delta p}
+
c^2[p]\partial_n\delta_\Gamma.
\end{equation}

Hence the ARD correction has the structure of a double-layer source
supported on the artificial interface.

\subsection{One-dimensional discretization}

For homogeneous Neumann conditions, it is natural to place the interface
between two grid nodes,
\[
\Gamma=x_{m+\frac12}.
\]
The continuous equation dictates that the correction is a dipole source:
\begin{equation}
\bar p_{tt}
=
c^2\overline{\Delta p}
+
c^2[p]\partial_n\delta_\Gamma.
\label{eq:wave-neumann-1d}
\end{equation}

To discretize this relation, consider the local state vector
\[
\mathbf p
=
\begin{pmatrix}
p_{m-1}\\
p_m\\
p_{m+1}\\
p_{m+2}
\end{pmatrix}.
\]

The pressure jump is approximated by
\begin{equation}
[p]
\approx
p_{m+1}-p_m.
\end{equation}

The interface lies halfway between the primary nodes,
\[
\Gamma=x_{m+\frac12}.
\]
Therefore the surface delta is naturally represented on the staggered
(face) grid by
\begin{equation}
\delta_\Gamma
\equiv
\frac1h
\begin{pmatrix}
0\\
1\\
0
\end{pmatrix}_{\mathrm{faces}},
\end{equation}
that is, a point source concentrated at the interface face
$x_{m+\frac12}$.

Taking the discrete normal derivative amounts to applying the discrete
divergence operator from the face grid to the primary grid. Consequently,
the derivative of the Dirac distribution is represented by the dipole
\begin{equation}
\partial_n\delta_\Gamma
\equiv
\frac1h
\begin{pmatrix}
-1 & 1
\end{pmatrix}
\left(
\frac1h
\begin{pmatrix}
0\\
1\\
0
\end{pmatrix}
\right)
=
\frac1{h^2}
\begin{pmatrix}
0\\
1\\
-1\\
0
\end{pmatrix},
\end{equation}
which is supported on the two grid nodes adjacent to the interface.

Therefore,
\begin{equation}
[p]\partial_n\delta_\Gamma
\approx
\frac1{h^2}
(p_{m+1}-p_m)
\begin{pmatrix}
0\\
1\\
-1\\
0
\end{pmatrix}
=
\frac1{h^2}
\begin{pmatrix}
0\\
p_{m+1}-p_m\\
-p_{m+1}+p_m\\
0
\end{pmatrix}.
\end{equation}

This vector can be written as a matrix acting on the state vector:
\begin{equation}
[p]\partial_n\delta_\Gamma
\approx
\frac1{h^2}
\begin{pmatrix}
0 & 0 & 0 & 0\\
0 & -1 & 1 & 0\\
0 & 1 & -1 & 0\\
0 & 0 & 0 & 0
\end{pmatrix}
\begin{pmatrix}
p_{m-1}\\
p_m\\
p_{m+1}\\
p_{m+2}
\end{pmatrix}.
\label{eq:discrete-neumann-matrix}
\end{equation}

\paragraph{Verification through the stencil defect.}

The global second-order Laplacian is
\begin{equation}
L_{\mathrm{global}}
=
\begin{pmatrix}
-2 & 1 & 0 & 0\\
1 & -2 & 1 & 0\\
0 & 1 & -2 & 1\\
0 & 0 & 1 & -2
\end{pmatrix}.
\end{equation}

Imposing homogeneous Neumann conditions at the interface yields the local
operator
\begin{equation}
L_N
=
\begin{pmatrix}
-2 & 1 & 0 & 0\\
1 & -1 & 0 & 0\\
0 & 0 & -1 & 1\\
0 & 0 & 1 & -2
\end{pmatrix}.
\end{equation}

Subtracting the local stencil from the global stencil gives
\begin{equation}
C_N
=
L_{\mathrm{global}}
-
L_N
=
\begin{pmatrix}
0 & 0 & 0 & 0\\
0 & -1 & 1 & 0\\
0 & 1 & -1 & 0\\
0 & 0 & 0 & 0
\end{pmatrix}.
\end{equation}

Applying this matrix to $\mathbf p$ gives
\[
C_N\mathbf p
=
(p_{m+1}-p_m)
\begin{pmatrix}
0\\
1\\
-1\\
0
\end{pmatrix},
\]
which coincides exactly with the discrete representation
\eqref{eq:discrete-neumann-matrix}. Therefore,
\begin{equation}
\frac1{h^2}C_N\mathbf p
\approx
[p]\partial_n\delta_\Gamma.
\end{equation}

Hence the ARD correction matrix is precisely the finite-difference
realization of the double-layer distribution generated by the artificial
Neumann interface.

\begin{thebibliography}{99}

\bibitem{Eldredge2021}
J.~D. Eldredge,
``A Method of Immersed Layers on Cartesian Grids, with Application to Incompressible Flows,''
\emph{Journal of Computational Physics},
vol.~448,
Art.~110716,
2022.
\href{https://doi.org/10.1016/j.jcp.2021.110716}
{doi:10.1016/j.jcp.2021.110716}.

\end{thebibliography}

\end{document}
